% TOP_PROBLEM: {"id": 92, "title": "Large Sum-Free Sets", "category_id": 2, "category": {"id": 2, "name": "combinatorics", "display_name": "Combinatorics", "description": "Counting problems, graph theory, discrete structures.", "slug": "combinatorics", "order_index": 2, "created_at": "2026-07-31T15:26:25.671Z"}, "status": "open"}
% FIRST_POSTED: null
% ENTRY_KIND: statement_only
% Imported from: batch02.tex; UnsolvedMath ID 92
\documentclass[11pt]{article}
\usepackage[margin=25mm]{geometry}
\usepackage{fontspec}
\setmainfont{Latin Modern Roman}
\usepackage{amsmath,amssymb,mathtools}
\usepackage{unicode-math}
\setmathfont{Latin Modern Math}
\usepackage{xurl,hyperref}
\hypersetup{colorlinks=true,urlcolor=blue}
\setcounter{secnumdepth}{0}
\setlength{\parindent}{0pt}
\setlength{\parskip}{5pt}
\setlength{\emergencystretch}{3em}
\newcommand{\sref}[2]{\hyperlink{source:#1:#2}{[S#2]}}
\newcommand{\cataloguescope}[1]{\paragraph{Recorded target.}#1}
\begin{document}
% UnsolvedMath ID 92; source catalogue code GREEN-001
\setcounter{section}{91}
\section{Large Sum-Free Sets}\label{92}
{\small\sffamily UnsolvedMath ID 92 · Combinatorics}
\subsection{Definitions and mathematical statement}

\paragraph{Shared notation.}$\mathbb N=\{1,2,\ldots\}$, and $\mathbb Z,\mathbb Q,\mathbb R,\mathbb C$ denote the integers, rationals, reals and complex numbers. Any other conventions are specified locally.


A subset $B\subset\mathbb Z$ is \emph{sum-free} if there are no $x,y,z\in B$ with $x+y=z$; the summands $x,y$ need not be distinct. For each finite $A\subset\mathbb N$, write
\[
s(A)=\max\{|B|:B\subset A\text{ is sum-free}\},\qquad
s(n)=\min_{\substack{A\subset\mathbb N\\|A|=n}}s(A).
\]
The catalogue's symbol for the secondary term denotes a function tending to infinity, rather than the usual asymptotic lower-bound notation.
\paragraph{Statement recorded in the supplied source.}
Does there exist a function $\omega:\mathbb N\to\mathbb R$ with $\omega(n)\to+\infty$ as $n\to\infty$ such that every $n$-element set $A\subset\mathbb N$ contains a sum-free subset $B\subset A$ satisfying
\[
|B|\ge \frac n3+\omega(n)
\]
for all sufficiently large $n$? Equivalently, is $s(n)-n/3\to+\infty$? \sref{92}{1}\sref{92}{2}
Bedert's Theorem 1.2 answers this historical question affirmatively by proving an additive improvement $c\log\log n$ for some absolute $c>0$. Here $\log$ denotes the natural logarithm, and the estimate is interpreted for sufficiently large $n$. \sref{92}{2}
\subsection{Short English statement}
Must every set of $n$ positive integers contain a sum-free subset whose size exceeds $n/3$ by an amount tending to infinity with $n$? Bedert's 2025 theorem gives an affirmative answer.
\subsection{Sources}
\begin{description}
\item[\hypertarget{source:92:1}{S1}] ULAM.AI and the UnsolvedMath Contributors, \emph{UnsolvedMath: A Curated Collection of Open Mathematics Problems}, record 92 (GREEN-001). CC BY 4.0. \url{https://huggingface.co/datasets/ulamai/UnsolvedMath}.
\item[\hypertarget{source:92:2}{S2}] Benjamin Bedert, \emph{Large sum-free subsets of sets of integers via $L^1$-estimates for trigonometric series}, arXiv:2502.08624v1, 12 February 2025, section 1, definition (1), Problem 1.1 and Theorem 1.2, pp. 1--2; Problem 1.1 explicitly identifies Green's Problem 1. \url{https://arxiv.org/pdf/2502.08624v1}.
\item[\hypertarget{source:92:3}{S3}] Ben Green, \emph{100 Open Problems}, author manuscript with December 2025 updates, Problem 1, pp. 2--3. \url{https://people.maths.ox.ac.uk/greenbj/papers/open-problems.pdf}.
\end{description}
\label{end:92}
\end{document}
